% !TEX root = ../../Monografia.tex
\chapter{Reprodutibilidade}
\label{ap:reprodutibilidade}

Todo o trabalho foi feito com código próprio, em Python, reunido no repositório do projeto
(\texttt{NewsManager}), a partir de dados públicos: o conjunto FNSPID \citep{dong2024fnspid} e os registros
da SEC. Este apêndice lista como cada resultado do texto é obtido.

\section{Etapas e comandos}

A Tabela~\ref{tab:comandos} lista as etapas na ordem de execução. Cada etapa grava seu produto em disco e
pode ser retomada.

\begin{table}[htbp]
\centering
\small
\caption{Etapas do processo e comando de cada uma.}
\label{tab:comandos}
\begin{tabularx}{\textwidth}{>{\raggedright\arraybackslash}p{4.0cm}>{\raggedright\arraybackslash}X}
\toprule
etapa & comando \\
\midrule
ingestão e curadoria do corpus & \texttt{python -m newsmanager \{probe, plan, ingest, curate, validate\}} \\
unidades e extração pelo LLM & \texttt{python -m newsmanager.extract \{units, run, collect, filter\}} \\
ligação a \emph{tickers} & \texttt{python -m newsmanager.extract \{link-prep, link-run, link-verify, link-build\}} \\
grafo de eventos & \texttt{python -m newsmanager.extract graph-build} \\
preços, universos e setores & \texttt{python -m newsmanager \{prices, universe, sectors\}} \\
regressão preditiva (E0--E4) & \texttt{python -m newsmanager h1} e \texttt{h1 --residual} \\
desenhos D1 e D2 & \texttt{python -m newsmanager backtest --n \{50, 100, 200\}} \\
desenhos D3 a D9 & \texttt{python -m newsmanager.backtester run backtests/*.toml} e \texttt{compare} \\
testes de pares (H1a, H1b) & \texttt{python -m newsmanager.grafo \{diadico-painel, diadico\}} \\
teste de permutação (QAP) & \texttt{scripts/revisao\_qap.py} \\
validação da extração e descritivas & \texttt{python -m newsmanager.grafo \{auditoria, amostra, metricas, anonimizacao, anon-run, anon-comparar, descritivas\}} \\
conferência dos números do texto & \texttt{scripts/conferir\_numeros.py} \\
\bottomrule
\end{tabularx}
\end{table}

\section{Versões e sementes}

\begin{itemize}
  \item Extração: versão \texttt{a05d0c0dc189}, \emph{hash} do prompt, do esquema, do limite de caracteres
        e do nome do modelo (Qwen2.5-14B-Instruct, servido pelo Ollama).
  \item Ligação a \emph{tickers}: versão \texttt{c49291db2953}, \emph{hash} do prompt, do modelo
        (Qwen2.5-7B-Instruct), dos tipos de agente e das listagens de referência.
  \item Parâmetros fixos: janela de notícias de 365 dias, meia-vida de 90 dias, janela de estimação de
        252 pregões, custo de 10 pontos-base, rebalanceamento no primeiro pregão do mês.
  \item Sementes: 7 nos \emph{bootstraps} e nas permutações do placebo; 101 a 104 no QAP (250
        permutações por semente).
  \item Cada desenho de carteira foi registrado no repositório antes da execução; o histórico de versões
        guarda o texto de cada pré-registro e a data.
\end{itemize}

\clearpage
